\chapter{Methodology}
\label{chap:methodology}
Chapter~\ref{chap:background} introduced the general tools this work relies
on. Here they are combined into a method that takes a case still in
execution and recommends which activity to perform next and which resource
should perform it, looking at several KPIs together. The KPIs often pull in
different directions, and the method does not settle the trade-off on behalf
of the decision-maker. It returns a small set of alternatives that balance
the KPIs in different ways.

The method has two parts. The first is built once, offline, from a log of
completed cases: one predictive model per KPI, each returning an estimate of
its KPI after a given next step together with a measure of how uncertain
that estimate is. The second part runs on every running case. It lists the
next steps that are plausible given the recent history of the case, scores
them with the models, discards the ones that are unlikely
to do better than letting the case proceed without a recommendation, and
picks a few representative steps from the Pareto front of those that remain.

The method uses only the event log and the trained models. The simulator
that estimates what would have happened had a recommendation been followed
plays no part in producing it: it serves to evaluate the recommendations and
is described in Chapter~\ref{chap:evaluation}.

% NOTE: the introduction overlaps with Section~\ref{sec:met:overview}; once
% that section is final, cut the introduction to two or three linking
% sentences.


\section{Overview of the Method}
\label{sec:met:overview}
The problem addressed in this thesis belongs to prescriptive process
analytics (Section~\ref{sec:bg:prpa}). While a case is running, the user
wants to know which next step to take, that is, which activity to perform
and which resource should perform it, so that the case ends as well as
possible with respect to the KPIs the user cares about. The method is split
into an offline and an online phase, shown in
Figure~\ref{fig:met:overview}.

The offline phase runs once, on a log of completed cases. The log is first
preprocessed; when it records only the completion time of each activity,
this step also estimates when the activity started. Two things are then
built from it. One is the feasibility prescription module, which records
the activities observed to follow a given recent history and the resources
observed to perform them, and is used later to decide which next steps are
feasible. The other is a training set of prefixes, encoded
with a frequency encoder, on which a predictive model is trained for each
KPI. Each model is used as an ensemble, so every prediction comes with an estimate of its uncertainty,
and is calibrated so that its predictions and uncertainty estimates agree
with the errors it makes on data not used for training.

Online, the method works on one running case at a time. Its prefix is
encoded in the same way as the training prefixes. The feasibility
prescription module gives the next steps that are feasible from the current
state of the case,
and the models estimate the KPIs the case would reach after each of them. A
confidence filter keeps only the steps that, given the uncertainty of the
predictions, are likely enough to improve at least one KPI over letting the
case continue without a recommendation. The steps that pass are compared
through their Pareto front, and a selection rule picks a few of them spread
across the front, which are returned to the user as recommendations.

Consider, for instance, a loan application in which the bank has just sent
an offer to the customer, with two KPIs: the probability that the customer
accepts, to be maximized, and the remaining time of the case, to be
minimized. The feasible next steps may include calling the customer or
preparing a new offer, and each can be assigned to several employees. A call
made by one employee may raise the chance of acceptance but keep the case
open longer, while another step may close the case sooner with a lower
chance of acceptance. Neither is better on both KPIs, so both can lie on the
Pareto front and the user receives both. A step whose predicted improvement
is small compared with the uncertainty of the models is removed by the
confidence filter before the front is built.

\begin{figure}[h]
  \centering
  \includegraphics[width=\linewidth]{chapters/3_Methodology/3_images/method_diagram.drawio.png}
  \caption{Overview of the method. In the offline phase, the completed cases
    of the event log are used to build the feasibility prescription module,
    based on a transition system, and to train a
    regressor and a classifier on an encoded set of prefixes; both models are
    calibrated and used as ensembles. In the online phase, a running
    execution is encoded and passed to the recommendation module, which
    evaluates its feasible next steps; the confidence filter, the Pareto
    front and the selection rule reduce them to the final recommendations.}
  \label{fig:met:overview}
\end{figure}

% --- Outline (metodologia.txt) ---
% 3.1 Panoramica del metodo
% Una figura della pipeline:
% parte offline: log -> preparazione dati -> due modelli con incertezza;
% parte online: caso in corso -> candidati -> valutazione -> filtro di
% confidenza -> fronte di Pareto -> top-k.
% Poche righe di testo. Il grosso lo fa la figura. Una spiegazione con magari
% un esempio di quello che si vuole andare a fare senza entrare troppo nel
% dettaglio.


\section{Problem Definition}
\label{sec:met:problem}
A KPI is modelled as a function
$F : \mathcal{E}^{*} \times \mathbb{N} \to \mathbb{R}$ that, given a trace
$\sigma$ and a position $1 \le i \le |\sigma|$, returns the value
$F(\sigma, i)$ of the indicator for $\sigma$ as seen after its first $i$
events~\cite{le2025balancing}. The value may depend on the whole trace,
events after position $i$ included, while the position fixes the point from
which it is measured, so a KPI can be referred to a case that is still
running. The process is judged on $m$ KPIs $F_1, \dots, F_m$, each with a
direction: it is either to be maximized or to be minimized.

The prescription is a \emph{next step}, a pair
$(a, r) \in \mathcal{A} \times \mathcal{R}$ made of the activity to perform
next and the resource assigned to it. Given a running execution
$\sigma' = \langle e_1, \dots, e_k \rangle$, the problem is to choose the
step to perform as the event $e_{k+1}$ so that the completed trace attains
the best values of the KPIs. With a single KPI the best step is the one with
the best value. With several KPIs, they are not merged into a single value:
each is predicted and compared on its own, and the trade-off between them is
left to the user.


\section{Offline Phase}
\label{sec:met:offline}
The offline phase works on a log $\mathcal{L}$ of completed traces and runs
once, before any recommendation is requested. What it builds stays fixed
during the online phase, which only queries it.

\subsection{Feasibility Prescription Module}
\label{subsec:met:offline:feasible}
A running case cannot be followed by any activity performed by any resource.
Some activities only make sense after others, and each activity is carried
out by a limited group of people. A recommender that ignored these
constraints could suggest steps that are impossible in practice, and it
would ask the predictive models to score combinations that never occur in
the data. Following~\cite{le2025balancing}, the next steps that can be
recommended for a prefix are restricted to those supported by the
historical log. The component that does so is the \emph{feasibility
prescription module}, built offline from $\mathcal{L}$. It combines a
transition system, which gives the activities that can follow the prefix, a
resource filter, which gives the resources that can perform each of them,
and the exclusion of the activities that cannot be prescribed.

The candidate activities come from a transition system mined from
$\mathcal{L}$ (Section~\ref{subsec:bg:pm:ts}), whose states keep the last
$w$ activities of a prefix:
\begin{equation}
  S^{w}(\sigma_{(k)}) =
  \begin{cases}
    \langle act(e_{k-w+1}), \dots, act(e_k) \rangle & \text{if } k \ge w,\\
    \langle act(e_1), \dots, act(e_k) \rangle & \text{otherwise.}
  \end{cases}
  \label{eq:met:off:state}
\end{equation}
The feasible next activities of a prefix $\sigma_{(k)}$ are those observed
in $\mathcal{L}$ right after its state:
\begin{equation}
  \mathcal{A}|_{\sigma_{(k)}} = \bigl\{\, a \in \mathcal{A} \bigm| \exists\,
  \tilde\sigma \in \mathcal{L},\ \exists\, l < |\tilde\sigma| :
  S^{w}(\tilde\sigma_{(l)}) = S^{w}(\sigma_{(k)}),\ a = act(\tilde e_{l+1})
  \,\bigr\},
  \label{eq:met:off:feasible}
\end{equation}
where $\tilde e_{l+1}$ is the $(l+1)$-th event of $\tilde\sigma$. An
activity is feasible if some historical trace reached the same state and
continued with it. A running case whose prefix never occurs in the log
still has feasible activities, as long as some other prefix reached the
same state. If the state itself was never reached in $\mathcal{L}$, the set
$\mathcal{A}|_{\sigma_{(k)}}$ is empty and no recommendation is produced for
the case.

The resource filter relies on the log to show who has performed what. An
activity $a$ can be assigned to any resource that performed it at least
once in $\mathcal{L}$:
\begin{equation}
  \mathcal{R}_a = \{\, r \in \mathcal{R} \mid \exists\, \sigma \in
  \mathcal{L},\ \exists\, e \in \sigma : act(e) = a,\ res(e) = r \,\}.
  \label{eq:met:off:resources}
\end{equation}
Unlike the activities, the eligible resources do not depend on the state of
the case.
% NOTE: the frequency thresholds (min_next_share, min_res_share) are
% described in the Evaluation chapter; there, note that with min_res_share the
% eligible resources depend on the window of the case, unlike in the method.

A last restriction concerns activities whose occurrence alone fixes the
value of a KPI. If a positive outcome is defined as the occurrence of an
activity, for instance the one recording that a customer accepted a loan
offer, recommending it would meet the KPI by definition without changing
anything in the case: accepting the offer is the customer's decision and
cannot be prescribed. An activity that marks a negative outcome is excluded
for a similar reason, since its effect on the KPI is known in advance and
there is nothing left for the models to estimate. These activities form a
set $\mathcal{A}_{F} \subseteq \mathcal{A}$, fixed for each process together
with its KPIs. The space of feasible next steps for $\sigma_{(k)}$ is then
\begin{equation}
  \mathcal{P}_{\sigma_{(k)}} = \{\, (a, r) \mid a \in
  \mathcal{A}|_{\sigma_{(k)}} \setminus \mathcal{A}_{F},\ r \in \mathcal{R}_a
  \,\} \subseteq \mathcal{A} \times \mathcal{R},
  \label{eq:met:off:steps}
\end{equation}
and these are the candidates evaluated in the online phase.

% --- Outline (metodologia.txt) ---
% 3.2 fase offline
% parlare del modello in generale e qual e' il suo obiettivo.
% definizione del transition system o citazione di quest'ultimo, a cosa serve
% e perche' lo usiamo (non dire che mettiamo una threshold perche' ne parlero'
% successivamente)
% [Il transition system con finestra w (le ultime w attivita') e l'insieme
% delle attivita' ammissibili.
% I due filtri sulla frequenza:
% min_next_share: si scartano le attivita' rare dopo quella finestra;
% min_res_share: si tengono solo le risorse frequenti per quella attivita'
% dopo quella finestra, e la piu' frequente resta sempre.
% Le attivita' proibite. Vanno spiegate: raccomandare l'attivita' che
% definisce l'outcome sarebbe banale.
% Il risultato e' l'insieme di coppie (a, r) da valutare.]

\subsection{Training Set and Encoding}
\label{subsec:met:offline:training}
The predictive models learn from the prefixes of the completed traces in
$\mathcal{L}$, each paired with the next step that actually followed it. For
a trace $\sigma = \langle e_1, \dots, e_n \rangle$ and a position
$1 \le k \le n$, the next step is $(a_{k+1}, r_{k+1}) = (act(e_{k+1}),
res(e_{k+1}))$. For $k = n$ no event follows, and the pair is replaced by an
end marker, so the models also see cases that have just finished. Given $m$
KPIs $F_1, \dots, F_m$ (Section~\ref{sec:met:problem}), the targets of
the instance are $F_1(\sigma, k), \dots, F_m(\sigma, k)$, the values the KPIs
take on the whole trace as seen after its first $k$ events.
% NOTE: the background defines the k-prefix with k < |sigma|; here (as in the
% code of le2025balancing) training also uses k = |sigma| with the end marker.

Prefixes have different lengths, while the models expect a fixed number of
features. Each prefix is therefore encoded with \emph{activity frequency
encoding}, a vector with one entry per activity of $\mathcal{A}$ that counts
how many times the activity occurred before the last event $e_k$. The
activity and the resource of $e_k$ are kept as separate features, together
with a few temporal ones: the time elapsed since the case started, the time
since the previous event, the duration of $e_k$ and the day of the week on
which it started. The case and event attributes recorded in the log are
added as they are. The next activity and the next resource complete the
vector as two categorical features. Writing $\mathrm{enc}(\sigma_{(k)}, a,
r)$ for the feature vector of prefix $\sigma_{(k)}$ followed by step
$(a, r)$, the training set is
\begin{equation}
  \mathcal{D} = \bigl\{ \bigl( \mathrm{enc}(\sigma_{(k)}, a_{k+1}, r_{k+1}),\
  F_1(\sigma, k), \dots, F_m(\sigma, k) \bigr) \bigm| \sigma \in \mathcal{L},\
  1 \le k \le |\sigma| \bigr\}.
  \label{eq:met:off:trainset}
\end{equation}
Figure~\ref{fig:met:off:instance} shows the structure of one instance. The
same encoding is applied, in the online phase, to the prefix of a running
case.

\begin{figure}[h]
  \centering
  \providecolor{Darkcandyapplered}{rgb}{0.64, 0.0, 0.0}
  \begin{tikzpicture}[
      box/.style={draw, rounded corners=2pt, align=center, minimum height=2cm,
                  text width=1.95cm, font=\scriptsize, inner sep=2pt},
      node distance=0.1cm]
    \node[box, fill=gray!15] (cur) {\textbf{Last event}\\activity and\\resource of $e_k$,\\weekday};
    \node[box, fill=gray!15, right=of cur] (hist) {\textbf{Prefix history}\\one counter per\\activity of $\mathcal{A}$};
    \node[box, fill=gray!15, right=of hist] (time) {\textbf{Time}\\since start,\\since previous\\event, duration\\of $e_k$};
    \node[box, fill=gray!15, right=of time] (attr) {\textbf{Attributes}\\case and event\\attributes\\of the log};
    \node[box, fill=Darkcandyapplered!15, right=0.3cm of attr, text width=2.1cm] (next) {\textbf{Next step}\\$(a_{k+1}, r_{k+1})$};
    \node[box, fill=Darkcandyapplered!75, text=white, right=0.5cm of next, text width=1.9cm] (y) {\textbf{Targets}\\$\mathrm{KPI}_1$\\$\vdots$\\$\mathrm{KPI}_m$};
    \draw[-{Latex}] (next) -- (y);
    \draw[decorate,decoration={brace,amplitude=5pt,mirror}]
      ([yshift=-3pt]cur.south west) -- ([yshift=-3pt]attr.south east)
      node[midway,below=6pt,font=\scriptsize] {state after $e_1, \dots, e_k$};
    \node[below=0.15cm of next, font=\scriptsize, align=center]
      {end marker after\\the last event;\\$\varnothing$ in $\mathcal{D}^{\varnothing}$};
  \end{tikzpicture}
  \caption{Structure of a training instance. The first four groups of
    features encode the state of the case after its first $k$ events. The
    next step is the one that followed $e_k$ in the log, replaced by an end
    marker after the last event of the trace and by $\varnothing$ in the
    copies of $\mathcal{D}^{\varnothing}$. The targets are the values of the
    $m$ KPIs.}
  \label{fig:met:off:instance}
\end{figure}

Every instance of $\mathcal{D}$ is then duplicated. The copy keeps the
prefix and the targets, and replaces the next activity and the next resource
with a token $\varnothing$ meaning that no next step is given:
\begin{equation}
  \mathcal{D}^{\varnothing} = \bigl\{ \bigl( \mathrm{enc}(\sigma_{(k)},
  \varnothing, \varnothing),\ F_1(\sigma, k), \dots, F_m(\sigma, k) \bigr)
  \bigm| \sigma \in \mathcal{L},\ 1 \le k \le |\sigma| \bigr\}.
  \label{eq:met:off:baselineset}
\end{equation}
The models are trained on $\mathcal{D} \cup \mathcal{D}^{\varnothing}$. The
online phase compares every candidate step with what would happen if the
case continued without a recommendation, and it needs a prediction for that
situation as well. A model trained only on instances that carry a next step
has never seen one without it, and its output for such an input would be an
extrapolation. An instance of $\mathcal{D}^{\varnothing}$ shows the
model a prefix without saying how the trace went on, together with the KPI
values that trace actually reached. Traces that pass through similar states
continue in different ways, so for the same input the model sees targets
produced by different next steps, and the best it can predict is their
average, weighted by how often each step was taken in the log. This average
is the value the case is expected to reach if it continues as similar cases
usually did, and it is the reference against which the candidate steps are
compared.

% --- Outline (metodologia.txt) ---
% encoding usato -> activity frequency encoder
% dire che addestriamo il modello dandogli sia la prossima attivita' che
% dandogli niente: spiegare quindi il perche'.

\subsection{Predictive Models}
\label{subsec:met:offline:models}
A separate predictive model is trained for each KPI, so no weight between
the KPIs is needed at training time, and each model can be of the kind its
KPI needs.
All the models take the same input, the vector
$\mathrm{enc}(\sigma_{(k)}, a, r)$, and are trained on
$\mathcal{D} \cup \mathcal{D}^{\varnothing}$.

The type of model depends on the KPI. A KPI with two possible values is
learned by a classifier, whose output is a logit $z$ that the logistic
function turns into a probability. A KPI with continuous values is learned
by a probabilistic regressor, which predicts a mean $\mu$ and a variance
$\sigma^{2}$ instead of a single value
(Section~\ref{subsec:bg:uncertainty:ensembles}). Both are trained by
maximizing the likelihood of the observed targets.
% NOTE (Evaluation): the remaining-time regressor is fitted in log(1 + days)
% space (chosen with 2b_compare_time_target.py), with the predictive
% log-normal truncated at the training maximum; predictions are mapped back
% to seconds.

Neither model is used as a single predictor. Each is used as an ensemble of
$M$ members (Section~\ref{subsec:bg:uncertainty:ensembles}), and every
member returns its own prediction for the same input: a logit $z_m$ for the
classifier, a pair $(\mu_m, \sigma_m^{2})$ for the regressor, with
$m = 1, \dots, M$. The predicted value $\hat{F}_j$ of KPI $F_j$ is the point
prediction of its model, a probability or a mean, and the members measure
how uncertain it is; for the regressor this uncertainty is the total
variance of Equation~\eqref{eq:bg:unc:gaussvar}, which adds the spread of
the means to the average of the variances. The online phase uses both
levels: the predicted values to place the candidate steps in objective
space, the members to judge whether a candidate is likely to do better than
the baseline.

Part of the training traces is held out as a validation set, split by trace
so that all the prefixes of a case, together with their copies in
$\mathcal{D}^{\varnothing}$, stay on the same side. The validation set is
used to select the hyperparameters of each model and to calibrate it
(Section~\ref{subsec:bg:uncertainty:calibration}). For the classifier, a
temperature $T$ is fitted as in Equation~\eqref{eq:bg:unc:tempfit} and every
logit, including those of the members, is divided by it before use. For the
regressor, a factor $s$ is fitted as in Equation~\eqref{eq:bg:unc:sigmafit}
and every predicted standard deviation is multiplied by it. Neither
correction changes the ranking of candidates by predicted value: they only
make the probabilities and the uncertainties that reach the online phase
agree with the errors observed on the held-out traces.
% NOTE (Evaluation): members are a virtual ensemble of a single SGLB model
% (truncations, Section bg:uncertainty:gbm); the classifier's point value is
% the calibrated probability of the full model, the regressor's the mean of
% the members. Validation = the last 20% of the training traces by completion
% time; hyperparameters by Optuna (minimizing the validation loss, early
% stopping); the final model is refitted on all of D u D^empty with the tree
% count of the best trial; T and s are fitted on that trial's model and reused
% on the final one (assumption: refitting on slightly more data does not
% change the miscalibration).

% --- Outline (metodologia.txt) ---
% modello che predice i kpi: regressore e classificatore (modello quindi piu'
% generico perche' usiamo kpi diversi e quindi devono essere trattati in modo
% diverso) scrivere anche il training e obiettivo
% c'e' anche da dire che di avere il risultato noi andiamo a fare anche la
% calibrazione del risultato del modello quando facciamo la exaustive search


\section{Online Phase}
\label{sec:met:online}
The online phase runs whenever a recommendation is requested for a running
case $\sigma_{(k)}$. It relies on the feasible steps and on the calibrated
models built offline, and needs nothing from the case beyond its prefix.

\subsection{Evaluating the Candidates}
\label{subsec:met:online:eval}
The prefix is encoded as in Section~\ref{subsec:met:offline:training}, and
the feasibility prescription module gives the space
$\mathcal{P}_{\sigma_{(k)}}$ of feasible next steps (Equation~\eqref{eq:met:off:steps}). Every pair
$(a, r)$ in it is a \emph{candidate}, denoted $C$, and its input to the
models is $x_C = \mathrm{enc}(\sigma_{(k)}, a, r)$. The \emph{baseline},
denoted $B$, is the same prefix with no next step,
$x_B = \mathrm{enc}(\sigma_{(k)}, \varnothing, \varnothing)$, which is the
kind of input the models learned from $\mathcal{D}^{\varnothing}$. A
candidate and the baseline differ only in the next activity and the next
resource, so any difference between their predictions is due to the step.

The candidates are evaluated by \emph{exhaustive search}: every pair in
$\mathcal{P}_{\sigma_{(k)}}$ is passed to every model, and none is skipped.
The result is collected by the function
\begin{equation}
  U_{\mathrm{KPI}} : \mathcal{P}_{\sigma_{(k)}} \to \mathbb{R}^{m}, \qquad
  U_{\mathrm{KPI}}(C) = \bigl( \hat{F}_1(C), \dots, \hat{F}_m(C) \bigr),
  \label{eq:met:on:objectives}
\end{equation}
which maps each candidate to the predicted values of the $m$ KPIs.
$\hat{F}_j(C)$ is the calibrated prediction of the model for KPI $F_j$ on
the input $x_C$ (Section~\ref{subsec:met:offline:models}), an estimate of
the value $F_j(\sigma, k)$ that the KPI would take on the completed trace
$\sigma$ if the case continued with the step $C$. Each KPI comes with a
direction, to be maximized or to be minimized, fixed together with the KPI
itself. The
baseline is passed to the models in the same batch. It is not a candidate
itself, but every candidate is compared with it.

Scoring every pair makes the search exact with respect to the models: the
best candidates according to the predictions are always found, and none is
lost because the search did not reach it. It requires
$\mathcal{P}_{\sigma_{(k)}}$ to be small enough to be scored in full.

% --- Outline (metodologia.txt) ---
% 3.3 fase online
% dato il prefisso e i modelli della fase online che valutano tutte le coppie
% noi applichiamo una exaustive search (spiegare bene cosa e') ovvero
% analizziamo tutte le possibili coppie e ne scegliamo quella migliore.
% dobbiamo quindi andare ad introdurre il concetto di metodo e di baseline.

\subsection{Probability of Improving on the Baseline}
\label{subsec:met:online:prob}
The predicted values say which candidate looks better than the baseline, but
not how far that judgment can be trusted. A candidate that beats the
baseline by less than the uncertainty of the models may owe its advantage to
noise in the predictions, and recommending it would mean intervening in the
case for no real gain. For every candidate and every KPI, the method
therefore estimates the probability that the candidate does better than the
baseline, using the members of the ensemble.

For KPI $F_j$, let $v_m(x)$ be the value that member $m$ assigns to the
input $x$, on the scale on which the comparison is made: the calibrated
logit $z_m / T$ for a classifier, the mean $\mu_m$ for a regressor. The
candidate and the baseline are not independent. Both are predicted by the
same members from the same prefix and differ only in the next step, so a
member that overestimates the KPI for this prefix tends to overestimate it
for both. Treating the two predictions as independent would count this
shared error twice and make every comparison look more uncertain than it
is. The comparison is therefore made on the paired differences
\begin{equation}
  d_m = v_m(x_C) - v_m(x_B), \qquad m = 1, \dots, M,
  \label{eq:met:on:diff}
\end{equation}
whose mean and sample variance are
\begin{equation}
  \bar{d} = \frac{1}{M} \sum_{m=1}^{M} d_m, \qquad
  \hat{\sigma}_d^{2} = \frac{1}{M-1} \sum_{m=1}^{M} \bigl( d_m - \bar{d}
  \bigr)^{2} = \hat{\sigma}_C^{2} + \hat{\sigma}_B^{2} -
  2\, \widehat{\mathrm{Cov}}(C, B),
  \label{eq:met:on:diffstats}
\end{equation}
where $\hat{\sigma}_C^{2}$ and $\hat{\sigma}_B^{2}$ are the sample variances
of the member values of the candidate and of the baseline, and
$\widehat{\mathrm{Cov}}(C, B)$ is their sample covariance. A positive
covariance, the usual case, removes from the comparison the part of the
uncertainty that the two predictions share, as in the example of
Figure~\ref{fig:met:on:paired}.

\begin{figure}[h]
  \centering
  \providecolor{Darkcandyapplered}{rgb}{0.64, 0.0, 0.0}
  \begin{tikzpicture}[font=\small]
    % (a) paired member values
    \begin{axis}[
        name=pairs, width=4.8cm, height=6.2cm,
        xmin=-0.35, xmax=1.35, ymin=-0.95, ymax=1.3,
        xtick={0,1}, xticklabels={$x_B$, $x_C$},
        ylabel={member value $v_m$}, ytick={-0.5,0,0.5,1},
        title={(a) paired members}, title style={font=\small},
        axis lines=left, tick style={black}]
      \foreach \b/\c in {0.08/0.41, -0.08/0.38, 0.38/0.36, 0.06/0.47,
                         -0.32/-0.12, 0.22/0.52, 0.78/1.12, 0.57/0.96,
                         -0.42/0.11, -0.76/-0.10} {
        \edef\temp{\noexpand\addplot[gray!60, mark=*, mark size=1.4pt,
          mark options={fill=gray!20, draw=gray!70}] coordinates {(0,\b) (1,\c)};}
        \temp
        \edef\temp{\noexpand\addplot[only marks, mark=*, mark size=1.4pt,
          Darkcandyapplered] coordinates {(1,\c)};}
        \temp
      }
    \end{axis}
    % (b) distribution of the difference
    \begin{axis}[
        at={(pairs.east)}, anchor=west, xshift=1.4cm,
        width=7.4cm, height=6.2cm,
        xmin=-1.6, xmax=2.8, ymin=0, ymax=2.35,
        xlabel={difference $d = v(x_C) - v(x_B)$}, ylabel={density},
        ytick=\empty, xtick={-1,0,0.36,1,2},
        xticklabels={$-1$,$0$,$\bar{d}$,$1$,$2$},
        title={(b) distribution of the difference}, title style={font=\small},
        axis lines=left, tick style={black}, samples=200, clip=false,
        declare function={
          tpdf(\x,\m,\s) = 0.38803/\s * (1 + ((\x-\m)/\s)^2/9)^(-5);}]
      \addplot[fill=Darkcandyapplered!25, draw=none, domain=0:2.8]
        {tpdf(x,0.36,0.194)} \closedcycle;
      \addplot[Darkcandyapplered, very thick, domain=-1.6:2.8]
        {tpdf(x,0.36,0.194)};
      \addplot[gray, thick, dashed, domain=-1.6:2.8]
        {tpdf(x,0.36,0.647)};
      \draw[black, thin] (axis cs:0,0) -- (axis cs:0,2.35);
      \node[Darkcandyapplered, anchor=north west, font=\footnotesize]
        at (axis cs:0.6,2.3) {paired: $p_j(C) = 0.95$};
      \node[gray!80!black, anchor=north west, font=\footnotesize]
        at (axis cs:0.6,2.05) {as if independent: $0.70$};
    \end{axis}
  \end{tikzpicture}
  \caption{Comparison of a candidate with the baseline on a synthetic
    example with $M = 10$ members, for a KPI to be maximized. (a) The value
    each member assigns to the baseline and to the candidate; a line joins
    the two values of the same member. The members disagree on the level of
    the KPI, but almost all of them place the candidate above the baseline.
    (b) The Student's $t$-distribution of the paired difference
    (Equation~\eqref{eq:met:on:prob}), whose shaded area to the right of zero
    is $p_j(C)$, and, dashed, the distribution obtained by ignoring the
    covariance and treating the two predictions as independent.}
  \label{fig:met:on:paired}
\end{figure}

For a regressor, the disagreement among the members covers only the
epistemic part of the uncertainty. The quantity being compared is the value
the KPI will take, so the aleatoric part is added as well, through the
average predicted variance of each input, and the result is rescaled by the
calibration factor $s$:
\begin{equation}
  \hat{\sigma}_d^{2} \leftarrow s^{2} \Bigl( \hat{\sigma}_d^{2} +
  \frac{1}{M} \sum_{m=1}^{M} \sigma_m^{2}(x_C) +
  \frac{1}{M} \sum_{m=1}^{M} \sigma_m^{2}(x_B) \Bigr).
  \label{eq:met:on:regvar}
\end{equation}
The aleatoric terms do not enter the covariance, since the noise that will
affect the case if the candidate is followed is not tied to the noise that
would affect it otherwise. For a classifier, only the disagreement among the
members is used. The variance of a single member would be the noise of the
binary outcome itself, while the quantity being compared is the probability,
and its uncertainty is the spread of the members. The comparison runs on
logits because a Gaussian fits an unbounded quantity better than one
confined to $[0, 1]$, and since the logistic function is monotone, comparing
logits and comparing probabilities give the same answer. Dividing every
logit by $T$ scales $\bar{d}$ and $\hat{\sigma}_d$ by the same factor, so the
result does not depend on the temperature.

Mean and variance are estimated from only $M$ members, and plugging them
into a Normal distribution as if they were exact would understate the
uncertainty of the comparison, the more so the smaller $M$ is. Given $M$
observations from a Normal population whose mean and variance are both
unknown, a new observation follows a Student's $t$-distribution with $M-1$
degrees of freedom, centered at the sample mean and with scale
$\hat{\sigma}_d \sqrt{1 + 1/M}$~\cite{gelman2013}. Applied to the paired
differences, this gives the probability that the candidate improves on the
baseline for KPI $F_j$:
\begin{equation}
  p_j(C) =
  \begin{cases}
    \mathcal{T}_{M-1}\!\left( \dfrac{\bar{d}}{\hat{\sigma}_d \sqrt{1 + 1/M}}
    \right) & \text{if } F_j \text{ is to be maximized},\\[3ex]
    1 - \mathcal{T}_{M-1}\!\left( \dfrac{\bar{d}}{\hat{\sigma}_d \sqrt{1 + 1/M}}
    \right) & \text{if } F_j \text{ is to be minimized},
  \end{cases}
  \label{eq:met:on:prob}
\end{equation}
where $\mathcal{T}_{M-1}$ is the cumulative distribution function of the standard
Student's $t$-distribution with $M-1$ degrees of freedom. As $M$ grows, the
$t$-distribution approaches the Normal and $\sqrt{1 + 1/M}$ approaches one.

% --- Outline (metodologia.txt) ---
% spiegare il filtro di confidenza senza andare a parlare del fatto che sto
% usando catboost ma comunque dal mio modello ottengo un certo ensemble e da li
% faccio le gaussiane se il numero e' sufficientemente alto eccetera. spiegare
% che facciamo la differenza di gaussiana tra baseline e metodo e che siccome
% sono dipendenti (spiega anche perche') devo fare anche covarianza (non
% chiamarlo a e b ma B per baseline e M per metodo -> C per candidate)
% [Per ogni candidato si stima P(il candidato batte la baseline senza
% raccomandazione) su ciascun KPI.
% Si usano gli M membri dell'ensemble, appaiati tra candidato e baseline:
% differenza gaussiana con la covarianza stimata;
% distribuzione t di Student con M-1 gradi di liberta';
% per classificazione binaria sulla scala dei logit;
% per regressione con la varianza totale per il sigma scale.]

\subsection{Confidence Filter}
\label{subsec:met:online:filter}
The probabilities of Equation~\eqref{eq:met:on:prob} are used to discard
the candidates that are not likely enough to improve on the baseline. For
every KPI a threshold $\gamma_j \in [0, 1]$ is fixed, and a candidate is
kept if it passes the threshold for at least one KPI:
\begin{equation}
  \mathcal{P}^{\gamma}_{\sigma_{(k)}} = \bigl\{\, C \in
  \mathcal{P}_{\sigma_{(k)}} \bigm| \exists\, j \in \{1, \dots, m\} :
  p_j(C) > \gamma_j \,\bigr\}.
  \label{eq:met:on:filter}
\end{equation}
A higher $\gamma_j$ makes the filter stricter on KPI $F_j$. Two candidates
with the same predicted values can therefore be treated differently: one is
kept because the members agree that it improves on the baseline, the other
is discarded because its advantage lies within their disagreement.

One KPI is enough to keep a candidate. A step that improves one KPI and
worsens another is still a valid option for the user, and requiring an
improvement on every KPI would discard it. The filter does not weigh the
gain on one KPI against the loss on another; it only checks that the gain is
credible.

If no candidate passes the filter, no recommendation is produced for the
case, since none of its feasible steps is likely to do better than letting
it continue without one.

% --- Outline (metodologia.txt) ---
% impostando un certo filtro gamma vado a togliere le raccomandazioni che per
% ogni kpi la probabilita' che il metodo sia migliore della baseline sia
% maggiore di gamma, rispettivamente 1 - prob nel caso in cui il kpi deve
% essere minimizzato.
% Regola: un candidato resta se P_kpi1 > gamma_1 oppure P_kpi2 > gamma_2 (OR)
% oppure P_kpi3.... Questo e' il punto in cui l'incertezza entra nella
% decisione, ed e' il contributo piu' originale del lavoro: gli darei spazio.

\subsection{Pareto Front of the Candidates}
\label{subsec:met:online:front}
The candidates that pass the filter are compared through their predicted
values $U_{\mathrm{KPI}}(C)$ (Equation~\eqref{eq:met:on:objectives}). The
definitions below are written for KPIs to be minimized, so that ``better''
always means ``lower''. A KPI to be maximized is brought to this form by
changing its sign, that is, by using $-\hat{F}_j$ in place of $\hat{F}_j$: a
higher value of the KPI becomes a lower value of $-\hat{F}_j$, and the
candidate preferred on that KPI stays the same. Following the notion of
\emph{Pareto optimality}~\cite{pareto1906,miettinen1999}, a candidate $C$
\emph{dominates} a candidate $C'$, written $C \prec C'$, when it is no worse
on every KPI and strictly better on at least one:
\begin{equation}
  C \prec C'
  \iff
  \bigl( \hat{F}_j(C) \le \hat{F}_j(C')\ \ \forall\, j \in
  \{1, \dots, m\} \bigr)
  \ \wedge\
  \bigl( \exists\, l : \hat{F}_l(C) < \hat{F}_l(C') \bigr).
  \label{eq:met:on:dominance}
\end{equation}
A candidate is \emph{Pareto-optimal}, or non-dominated, if no other
candidate dominates it, and the non-dominated candidates of the case are
\begin{equation}
  \mathcal{F}_{\sigma_{(k)}} = \bigl\{\, C \in
  \mathcal{P}^{\gamma}_{\sigma_{(k)}} \bigm| \nexists\, C' \in
  \mathcal{P}^{\gamma}_{\sigma_{(k)}} : C' \prec C \,\bigr\}.
  \label{eq:met:on:front}
\end{equation}
Their predicted values form the \emph{Pareto front} of the case, shown for
two KPIs in Figure~\ref{fig:met:on:pareto}. The point that has the best
possible value of every KPI is the \emph{ideal point}. It is generally not
attainable, because the KPIs conflict, and the front shows how close the
candidates can get to it.

\begin{figure}[h]
  \centering
  \providecolor{Darkcandyapplered}{rgb}{0.64, 0.0, 0.0}
  \begin{tikzpicture}[font=\small]
    \begin{axis}[
        width=9cm, height=7cm,
        xmin=0, xmax=8.4, ymin=0, ymax=7.4,
        axis lines=left, xtick=\empty, ytick=\empty,
        xlabel={$\hat{F}_1$}, ylabel={$\hat{F}_2$},
        xlabel style={at={(axis description cs:1,0)}, anchor=north east},
        ylabel style={at={(axis description cs:0,1)}, anchor=south east,
                      rotate=-90},
        clip=false,
        legend style={at={(1.03,0.98)}, anchor=north west, font=\footnotesize,
                      draw=gray!50, fill=white, row sep=1pt},
        legend cell align=left]
      % region dominated by p
      \fill[gray!15] (axis cs:2.4,3.6) rectangle (axis cs:8.4,7.4);
      \node[gray!70!black, font=\footnotesize\itshape] at (axis cs:5.9,6.55)
        {candidates dominated by $p$};
      % front
      \addplot[Darkcandyapplered, dashed, thick, forget plot]
        coordinates {(1.0,6.0) (1.6,4.6) (2.4,3.6) (3.4,2.9) (4.6,2.4)
                     (5.8,2.1) (7.2,1.9)};
      \addplot[only marks, mark=*, mark size=2.6pt, Darkcandyapplered]
        coordinates {(1.0,6.0) (1.6,4.6) (2.4,3.6) (3.4,2.9) (4.6,2.4)
                     (5.8,2.1) (7.2,1.9)};
      \addlegendentry{non-dominated candidate}
      % dominated
      \addplot[only marks, mark=o, mark size=2.6pt, gray, thick]
        coordinates {(2.8,5.6) (3.6,4.9) (5.0,4.4) (6.2,3.9) (4.3,3.4)
                     (7.0,3.3) (2.0,6.6) (6.4,2.8)};
      \addlegendentry{dominated candidate}
      % ideal point
      \addplot[only marks, mark=star, mark size=4pt, black, very thick]
        coordinates {(0.3,0.3)};
      \addlegendentry{ideal point}
      % ringed candidate p
      \draw[black, thick] (axis cs:2.4,3.6) circle[radius=4.2pt];
      \node[anchor=east] at (axis cs:2.25,3.45) {$p$};
      \node[Darkcandyapplered, font=\footnotesize, rotate=-17]
        at (axis cs:5.0,1.75) {Pareto front};
    \end{axis}
  \end{tikzpicture}
  \caption{Pareto front of a case with two KPIs, both to be minimized. Each
    marker is a candidate placed by its predicted values. An open circle is
    a dominated candidate (Equation~\eqref{eq:met:on:dominance}); the shaded
    area shows the candidates dominated by the ringed candidate $p$. The
    filled points are non-dominated and form the Pareto front, along which
    any gain on one KPI is paid for by a loss on the other. The ideal point,
    with the best possible value of both KPIs, lies at the bottom-left.}
  \label{fig:met:on:pareto}
\end{figure}

A dominated candidate is discarded because another feasible step is
predicted to be at least as good on every KPI and better on one, so the user
would have no reason to prefer it. The steps left on the front are those
among which a choice actually has to be made: moving from one to another
improves some KPI only at the cost of another, and none of them is better
than the others until a preference between the KPIs is stated. Returning the
front leaves that preference to the user, who chooses a step with the cost
of each alternative in view.

The front is computed on the predicted values as they are, without
normalizing them. Dominance compares the candidates one KPI at a time, so
rescaling a KPI would not change which candidates are dominated. The
baseline is not a candidate and does not take part in the comparison.

% --- Outline (metodologia.txt) ---
% una volta filtrati i candidati andiamo a costruire un fronte di pareto
% (definizione del fronte)
% Dominanza e fronte riscritti sul vostro problema: f(a, r). Va spostato dal
% background (par. 2.3.2).
% Si calcola il fronte sui valori non normalizzati.
% Ricerca esaustiva

\subsection{Selecting the Recommendations}
\label{subsec:met:online:selection}
The front of a case can contain many candidates, some of them very close to
one another in objective space. A user is better served by a few clearly
distinct options than by the whole front, so at most $k$ candidates are
returned. If $|\mathcal{F}_{\sigma_{(k)}}| \le k$, all of them are. Otherwise
the $k$ candidates are chosen to be as spread out as possible along the
front.

Choosing the $k$ most spread-out points of a set is the \emph{max-min
dispersion}, or \emph{$p$-dispersion}, problem: given $n$ points with
pairwise distances $d_{il}$, select a subset $S$ of size $k$ that maximizes
the smallest distance between any two selected points,
\begin{equation}
  \max_{\substack{S \subseteq \{1, \dots, n\} \\ |S| = k}}\ \
  \min_{\substack{i, l \in S \\ i \neq l}}\ d_{il}.
  \label{eq:met:on:dispersion}
\end{equation}
The problem was first studied in facility location, where the aim is to
place $p$ facilities so that the closest pair is as far apart as
possible~\cite{kuby1987}, and it has since been used more generally to
extract a spread-out subset from a set of candidates~\cite{maliszewski2012}.

Here the points are the candidates of $\mathcal{F}_{\sigma_{(k)}}$, placed
by their predicted values. The KPIs are measured on different scales, and a
distance computed on the raw values would be dominated by the KPI with the
widest range. Each objective is therefore rescaled to $[0, 1]$ between its
best and its worst value over $\mathcal{F}_{\sigma_{(k)}}$:
\begin{equation}
  \tilde{F}_j(C) = \frac{\hat{F}_j(C) - \min_{C' \in
  \mathcal{F}_{\sigma_{(k)}}} \hat{F}_j(C')}{\max_{C' \in
  \mathcal{F}_{\sigma_{(k)}}} \hat{F}_j(C') - \min_{C' \in
  \mathcal{F}_{\sigma_{(k)}}} \hat{F}_j(C')},
  \label{eq:met:on:minmax}
\end{equation}
with $\tilde{F}_j(C) = 0$ when the KPI takes the same value on the whole
front. The distance between two candidates is the Euclidean distance between
their rescaled values,
$d(C, C') = \lVert \tilde{F}(C) - \tilde{F}(C') \rVert_2$. The
rescaling serves only to measure distances: the front itself is computed on
the predicted values as they are.

Problem~\eqref{eq:met:on:dispersion} can be written as an integer linear
program and solved exactly when the number of points is small, as it is for
the front of a single case. The $k$ candidates it selects are returned as
the recommendations for the case.

% --- Outline (metodologia.txt) ---
% nel caso in cui i punti del fronte dovessero essere tanti si vuole solo
% andare ad avere un subset di questi e per farlo si usa il problema min-max
% delle distanze tra punti
% Normalizzazione min-max degli obiettivi, solo per calcolare le distanze.
% Problema max-min dispersion sulle distanze euclidee, restiamo generici. Va
% spostato dal background (par. 2.3.3).
% Se il fronte ha k punti o meno, si restituiscono tutti.
% L'output sono k coppie (a, r) per caso.
